\SourceLeaf{151}coherence lacking, and I see in it nothing more than a bare Petitio principii. You will not misunderstand me. I attach little importance to a strictly logical presentation, which in writings for men and experts would often be mere pedantry. A strict logical connection must, however, be demonstrable throughout whenever it is called for. That is an indispensable condition of a good exposition. If I think very highly of a mathematician and find something missing in this respect, I always think first that it is \textls{merely} in the presentation, and I am reluctant to assume that the idea itself is empty. In the case of Mr \textsc{Ivory}'s three articles I am unfortunately forced to do so.

The paper on least squares is indeed beneath all criticism. What confusion, obscurity, and complete lack of logical cogency! One could still overlook this if there were a genuine idea behind it. That is altogether lacking. How can the prattle on pp. 163 and 164 be passed off to us as a \textls{proof} that this principle leads \textls{necessarily} to the method of the least squares!

Far from showing the \textls{necessity} of the method of least squares, this whole argument\HistoricalNote{*)}{*)}{Except, of course, for the consequence that $\varepsilon\varepsilon+\varepsilon'\varepsilon'+\varepsilon''\varepsilon''+\cdots$ becomes a minimum. But this is precisely what is to be proved and must not be an \textls{axiom}.} can in fact be applied to \textls{any} other treatment of the given equations. For example, instead of multiplying them by $a$, $a'$, $a''$, one could multiply them by any power of these coefficients or even divide by them. The paper really seems to have been written half asleep. You will surely spare me pointing out several more individual passages, which you will notice at once yourself if you do not start from the preconceived opinion that there is anything in it.

The other paper on pendulum lengths seems to me little better. Its objections to applying the method of least squares show only that its spirit is entirely foreign to Mr I. I find that Mr \textsc{Sabine} applied this method quite correctly in essentials (— the form of the application is admittedly not the most expedient, and between ourselves I actually had \textsc{Sabine}, \textsc{Paucker}, and \textsc{Muncke} in the new edition of \textsc{Gehler} in mind in my
